\section{Cross-Domain Synthesis}
\label{sec:ieee-synthesis}

The scientific value of ORIUS is not that every domain already looks equally complete. It
is that one runtime grammar survives across very different physical systems while the
paper remains honest about the current evidence depth. That distinction
is what turns ORIUS from a portability story into a candidate field-level safety layer.

\begin{table*}[t]
\centering
\caption{What ORIUS standardizes across domains, and what remains domain-specific.}
\label{tab:ieee-universal-vs-domain}
\small
\begin{tabularx}{\textwidth}{@{}p{0.20\textwidth}p{0.32\textwidth}p{0.40\textwidth}@{}}
\toprule
\textbf{Layer} & \textbf{Shared across all rows} & \textbf{Domain-specific content}\\
\midrule
Observation semantics & Reliability assessment, degraded-observation detection, uncertainty widening trigger & Sensor modalities, freshness signals, and consistency features.\\
Safety geometry & Tightened admissible action set, repair-or-fallback release logic & State constraints, action constraints, safety-object geometry.\\
Benchmarking & True-state violation, intervention, fallback, certificate validity, latency, domain metrics & Domain-specific margin definitions and auxiliary metrics.\\
Governance & CertOS lifecycle, certificate integrity, audit continuity & Policy-specific fallback payloads and application-specific certificate fields.\\
\bottomrule
\end{tabularx}
\end{table*}

This synthesis yields a sharper universality statement than a transfer-learning claim.
ORIUS is not saying that one policy transfers across every domain. It is saying that
Physical AI systems share a runtime failure mode once observation degrades, and that the
same kernel can mediate that failure mode while respecting domain-specific safety
geometry. That is why the architecture can be universal even when the evidence remains
tiered.

The paper therefore keeps the evidence boundary visible. Battery is the primary
validation domain. AV is validated on all-zip grouped nuPlan replay. Healthcare is
validated as a bounded retrospective monitoring domain. Future additional domains are
not included in the public claim surface until they satisfy the same evidence criteria.
This is the mechanism that lets the paper make a large architectural claim without
pretending that universality of contract automatically implies equal evidence depth.

\subsection{Three-domain evidence boundary}
The public claim surface is deliberately three-domain. Battery establishes the primary
validation standard: theorem-facing assumptions, runtime traces, fallback semantics, and
CertOS evidence are all present in the same domain. The AV domain tests whether the same release
grammar survives a faster mobility setting, but its current denominator is bounded to
completed nuPlan replay artifacts and the TTC/predictive-entry-barrier contract. The
Healthcare domain tests the same grammar in retrospective monitoring, but it remains a
MIMIC-backed alert-release row and not a clinical deployment result. This boundary is
part of the contribution because it prevents architectural universality from being
misread as equal empirical maturity.

\paragraph{Headline synthesis.}
Across the three domains, ORIUS achieves zero or near-zero TSVR in every row:
$0.833\% \to 0.000\%$ (Battery), $28.93\% \to 0.016\%$ (AV, 1.53M steps),
and $19.45\% \to 0.000\%$ (Healthcare, 137K steps).
In every row, runtime latency remains sub-millisecond (p95: 0.81, 0.74, 0.70\,ms),
and certificate-valid release rate is above $99.3\%$.
These results are achieved without collapsing to degenerate fail-safes: ORIUS
retains 10.1\,MWh useful work over battery shutdown, $1.68\times$ the useful AV
driving work over always-brake, and 142,767 useful monitoring units over
always-alert healthcare, passing the utility-preserving safety gate in all three domains.
